\scanpage{12}

\deleted{$\unsure{I(U)}\otimes_k I(\Lambda)$ le soit.} D'où
\deleted{\diagram{\begin{tikzcd}
I(U)\otimes_k I(\Lambda) \arrow[r] & 1+I(U)\otimes_k
\end{tikzcd}}}
\inserted{$U\otimes_k I(\Lambda)$}
\figure{\begin{tikzpicture}\draw(0,1)..controls(1.2,1)and(.1,-.2)..(1,-.2)--(2.4,-.2)..controls(2.4,.5)and(2.6,.7)..(2.9,.7);\end{tikzpicture}}

Alors $U\otimes I(\Lambda)$ est de puiss. n-ième
nulle, d'où
\diagram{\begin{tikzcd}
U\otimes I(\Lambda) \arrow[r,bend left=15,"\exp"] & 1+U\otimes I(\Lambda) \arrow[l,bend left=15,"\log"]
\end{tikzcd}}
[et si $I(U)$ est l'idéal d'une
augmentation de $U$, on a
\diagram{\begin{tikzcd}
I(U)\otimes I(\Lambda) \arrow[r,bend left=15,"\exp"] & 1+I(U)\otimes I(\Lambda) \arrow[l,bend left=15,"\log"]
\end{tikzcd}}]

Supposons maintenant que $U$ soit une
hyperalgèbre ordinaire. Alors la correspondance
précédente établit une corr. biunivoque
entre la partie de
\deleted{$\unsure{I(U)\otimes I(\Lambda)}$}
\inserted{$(U\otimes I(\Lambda))\cap I(U\otimes\Lambda)
\unsure{\bigl(\Lambda\otimes_k I(U)\otimes I(\Lambda)\bigr)}$}
formée
des éléments primitifs \unsure{de} $U\otimes_k\Lambda$, et
la partie de $1+I(U)\otimes I(\Lambda)$
\deleted{$\unsure{\subset1+I(U\otimes\Lambda)}$}
formée des éléments multiplicatifs. Ce
sont aussi resp. les primitifs de
$U\otimes_k\Lambda$ qui donnent $0$ par l'augmentation
$U\otimes_k\Lambda\longrightarrow U$, et les éléments multiplicatifs de
$U\otimes_k\Lambda$ qui donnent $1$ par cette \unsure{augmentation}.

\mspage{11}
